\documentclass[../main.tex]{subfiles}
\begin{document}
%==========================================TIÊU ĐỀ TẬP========================================
\begin{table}[h]
  \begin{adjustwidth}{-1cm}{}
    \begin{tabular}{>{\centering\arraybackslash}p{6cm}|>{\raggedright\arraybackslash}p{12.5cm}}
      \multirow{5}{6cm}
      % Cột bên trái: ÉPISODE và số tập
      {\\[-40pt]
        \flaregothic\fontsize{20pt}{20pt}\selectfont \centering
        \color{eptype}ÉPISODE\color{black}\\
        \burbank\fontsize{72pt}{72pt}\selectfont
        \color{epnum}102
        \\[18pt]
        % \flaregothic\fontsize{14pt}{14pt}\selectfont
        % \color{epattr}BIENVENUE, \uline{\color{Banpire} Mel} !
        \color{black}
      }
       &
      % Dòng thứ nhất: tên tập tiếng Nhật
      {\fotskip\fontsize{18pt}{18pt}\selectfont まだ\ruby{誰}{だれ}にも\ruby{見}{み}せていない
      }   \\[6pt]
      % Dòng thứ hai: tên tập tiếng Anh
       & {\cafeteria\fontsize{18pt}{18pt}\selectfont The Realistic Look}          \\[4pt]
      % Dòng thứ ba: tên tập tiếng Việt
       & {\vnmsans\fontsize{18pt}{18pt}\selectfont Máy ảnh thực thể hóa}          \\[8pt]
      % Dòng thứ tư: Nhân vật xuất hiện
       & {\sen{\fontsize{14pt}{14pt}\selectfont Track used: Goblin Town (ゴブリンタウン)}
         }
      \\[8pt]
      % Dòng cuối cùng: Thời gian diễn ra của tập (thường sẽ là ngày phát hành)
       & {\continuum\fontsize{16pt}{16pt}\selectfont Ngày phát hành: 18/04/2021}  \\[18pt]
    \end{tabular}
  \end{adjustwidth}
\end{table}
%=========================================TRUYỆN CHÍNH========================================
\color{black}

Một buổi sáng nọ ở văn phòng. Mio, Mel và Kuon phát hiện một chiếc máy ảnh đặt trên bàn.

Nàng Sói đen dựng đôi tai sói của mình lên: ``Ồ, cái máy ảnh này của ai thế?''

Cậu Tổng nhún vai: ``Anh cũng chẳng biết. Lúc mới tới đã thấy nó ở đây rồi. Nhìn cũng na ná cái của anh ở nhà.''

Nàng Dơi cầm chiếc máy ảnh lên, săm soi một chút rồi nhíu mày tặc lưỡi: ``Năm 2021 rồi mà vẫn còn xài mấy cái máy này à?''

\img{8-6a}

Sói đen nhướng mày, đổ mồ hôi hột với suy nghĩ giản đơn của tiền bối: ``Ờm\dots\ chứ còn sao nữa?''

Cậu Tổng giải thích thêm: ``Vẫn dùng chứ, nhất là mấy tay thợ chụp ảnh chuyên nghiệp. Họ không giống như mấy đứa dùng điện thoại tự sướng rồi đăng lên mạng đâu.''

Mio nghe vậy, gật gù: ``Hơi phũ, nhưng công nhận là thế.''

``Nhưng quan trọng là nó còn dùng được hay không?'' Kuon nói, nhìn chiếc máy với vẻ nghi hoặc.

Mel búng tay, tỏ vẻ hứng thú: ``Hay là thử luôn đi!''

Vừa hay, họ thấy Watame đang nhảy nhót, hai tay đung đưa một cách nhiệt tình như đang tập thể dục buổi sáng. Không bỏ lỡ cơ hội, người đang cầm máy ảnh - là Mel - liền giơ máy lên, bấm chụp một cách nhanh gọn.

*TÁCH!*

Một luồng sáng kỳ lạ lóe lên từ ống kính. Cả ba nheo mắt theo phản xạ. Nhưng khi họ mở mắt ra\dots

``Ối chà!'' nàng Dơi thảng thốt với thứ mà cô nhìn thấy.

Mio toát mồ hôi hột: ``Watame\dots\ biến hình rồi kìa!''

Kuon nhìn cảnh tượng trước mắt, bật cười: ``Chuẩn bản chất thật của em ấy luôn!''

Đứng trước mặt họ không còn là Watame trong hình dạng idol, mà là một con cừu lông xù chính hiệu - hàng thật một trăm phần trăm!

\begin{figure}[H]
  \centering
  \begin{subfigure}[t]{0.45\textwidth}
    \centering
    \includegraphics[width=8cm]{../../../../Image/Book?/8-6b1.png}
    \caption{Watame trước khi chụp.}
  \end{subfigure}
  \quad
  \begin{subfigure}[t]{0.45\textwidth}
    \centering
    \includegraphics[width=8cm]{../../../../Image/Book?/8-6b2.png}
    \caption{Watame sau khi chụp.}
  \end{subfigure}
\end{figure}

Con cừu -tức Watame - hoảng loạn vẫy vùng, miệng phát ra những tiếng kêu be be hoảng hốt, như muốn hét lên: ``\textbf{\textit{Chuyện gì vậy!? Cứu tôi với!!}}''

Cả ba nhìn lại chiếc máy ảnh, cảm nhận thấy một thứ ma lực khủng khiếp tỏa ra từ nó.

Nàng Sói đen nuốt khan: ``Đáng sợ thật\dots''

Kuon thì rùng mình: ``Làm anh nhớ đến vụ camera của Ayame-chan gần hai năm trước\footnote{Xem lại {\flaregothic ÉPISODE 6}, quyển số Một.}\dots\ Hãi lắm.''

Nàng Ma cà rồng vẫn còn sốc, chỉ vào con cừu vẫn đang kêu lên hoảng loạn trước mặt: ``Đó\dots\ là Watame-chan sao?''

Cả Kuon và Mio đồng thanh: ``Thì em ấy còn ai! Tại cái camera đấy!''

Lo rằng cái máy ảnh này còn có thể gây ra những chuyện kinh khủng hơn, Kuon và Mio lập tức bảo Mel đặt nó lại chỗ cũ. Nhưng\dots

\dots Tri giác của nàng Dơi lại đang mách bảo làm gì đó với cái máy ảnh mới mẻ này. Một điều gì đó không được tốt lành.

Thay vì nghe lời, Mel lại cầm nó lên và phóng ra khỏi văn phòng.

Cậu Tổng thở dài xoa trán: ``Và\dots\ lại nước đổ đầu dơi rồi. Tò mò thì hại người mà.''

Nàng Sói đen khoanh tay, trừng mắt nhìn theo cánh cửa: ``Em cá là chị ấy sẽ lại nghịch cho mà xem! Lại tính mưu gì đây mà!''

Kuon chậc lưỡi: ``\dots Cũng giống vụ mà Roboco-chan cầm cái máy ảnh của Ayame-chan năm đó mà.''

Mio quay sang, lo sợ: ``Ý anh là sao?''

``Chuyện dài lắm. Hồi đó em còn chưa vào mà-'' cậu định kể tiếp, nhưng chợt nhớ ra nguy cơ trước mặt, bèn hối thúc: ``Ơ kìa! Đuổi theo em ấy đi, kẻo lại làm chuyện xấu thì phiền lắm!''

Ngay lập tức, cặp đôi ``thanh niên nghiêm túc'' vội vàng đuổi theo nàng Ma cà rồng, bỏ lại một con cừu đang nhảy cẫng lên giữa văn phòng, miệng la oai oái: ``\textbf{\textit{Quay tôi về đi mà!!!}}''

\ct

Mel cầm máy ảnh chạy tót ra ngoài đường, phấn khởi như một đứa trẻ mới tìm thấy đồ chơi mới. Cô cười tít mắt, lia máy ảnh khắp nơi, chụp lia lịa mà chẳng thèm để ý đến hậu quả.

Bất chợt, cô thấy Kanata và A-chan đang đứng gần đó, trò chuyện với nhau. Như một phản xạ tự nhiên, nàng Dơi vẫy tay gọi to: ``\textbf{A-chan! KanaKana!}''

Hai người quay lại theo phản xạ. nàng Thiên sứ hơi nghiêng đầu, vừa mở miệng chào: ``À, Mel-senp-''

*TÁCH!*

Chưa kịp nói hết câu, ánh sáng từ máy ảnh đã lóe lên, bao phủ cả một vùng sáng.

Nàng Đeo kính chói mắt, lập tức lấy tay che mặt:
``Ái da! Mắt của chị!!''

Khi ánh sáng dịu đi, cô bỏ tay ra, nhìn về phía nàng tóc vàng với ánh mắt trách móc: ``Chụp ảnh người khác mà không xin phép là bất lịch sự lắm nha!''

Cô lườm Mel một cái rồi - có lẽ để ``dạy cho một bài học'' - cũng giơ máy lên, bấm chụp một phát vào nàng Dơi đang cười trừ.

Nhưng\dots\ khi quay sang Kanata để tiếp tục càm ràm, cô chợt khựng lại. Cả khuôn mặt của nàng Dơi cũng xanh xao như tàu lá chuối.

Mặt A-chan dần tái mét.

Bởi trước mặt họ, tại vị trí Kanata đứng ban nãy\dots\ là một bức tượng cẩm thạch Thiên sứ biết nói. Đây là bức tượng với đôi cánh ngẩng lên trời, khuôn mặt nghiêm nghị với bộ trang phục trông như đến từ thời kỳ Hy Lạp cổ đại.

\img{8-6c}

Đôi mắt của nó sáng rực, tỏa ra một thứ ánh sáng thần thánh. Và rồi, giọng nói quen thuộc của Thiên sứ vang lên từ trong đó: ``\textbf{Hãy yêu quý những người bạn của ngươi.}''

Một luồng hào quang từ trên trời rọi xuống. A-chan sửng sốt, lùi lại một bước: ``Ôi thánh thần ơi!?''

Mel thì vẫn đứng chôn chân ở đó, không thốt ra một câu nói nào.

\ct

Trong khi đó, Kuon và Mio vẫn đang chạy khắp thị trấn để truy tìm nàng Dơi trốn chạy.

Nàng Sói chống đầu gối thở dốc: ``el-chan chạy đi đâu rồi!?''

Cậu Tổng vừa chạy vừa nhíu mày nhìn quanh: ``Không biết - À, em ấy ở kia kìa! A-san cũng ở đó nữa, và\dots\ oắt-đờ-gót!?''

Cả hai nhanh chóng chạy lại gần.

Trước mắt họ là một khung cảnh mà ngay cả trong giấc mơ cũng chưa từng nghĩ đến: Mel và A-chan đang đứng trước một bức tượng cẩm thạch khổng lồ\dots\ và bức tượng phát ra giọng nói của Kanata.

Một khoảng im lặng nặng nề bao trùm. Rồi cuối cùng, họ chỉ có thể bật ra một câu: ``C-Cái gì đây\dots!?''

Bức tượng (Kanata) từ trên cao nhìn xuống, giọng nói vang vọng đầy uy nghiêm: ``\textbf{Những cái xấu xa đã chiếm lấy ngươi rồi. Các ngươi sẽ cần được thanh tẩy từ ta đây!}''

Cậu Tổng chớp mắt, hoàn toàn không hiểu chuyện gì đang diễn ra: ``Cái gì cơ!?''

Nàng Sói đen trừng mắt, hét lên: ``Cái quái gì thế!?''

``Anh không theo đạo Thiên Chúa, anh không có biết đâu,'' cậu khoanh tay lắc đầu.

Bức tượng nhìn chằm chằm xuống họ, giọng nói trở nên vang vọng hơn: ``\textbf{Mong các ngươi sẽ được cứu rỗi khỏi cái chết!}''

Nói rồi, nó lóe lên một ánh màu vàng rực rỡ.

\img{8-6d}

Mio sững người: ``Khoan\dots\ không lẽ nó có quyền quyết định tối cao đến vậy sao!?''

Kuon bối rối lùi lại, quát lên: ``Làm sao mà anh biết được!? Với cả\dots''

Cả hai cùng hét lên: ``\textbf{Tại sao chúng ta cũng bị liên lụy chứ!??}''

Thấy tình hình có vẻ không ổn, Sói đen bảo với nàng Dơi rằng: ``Mel! Thử chụp cái bức tượng ấy xem nào!''

Nghe lệnh của cô, người cầm máy ảnh nói: ``OK!'' rồi nhanh chóng chụp nó, bức tượng thánh thần kia liền trở lại Kanata.

``Ủa? Mình đang làm gì ở đây vậy?''

Chưa kịp thở phào nhẹ nhõm, thì chẳng biết thế nào, Mel lại chụp Kanata lần nữa, biến cô nàng Thiên sứ trở lại bức tượng.

\begin{figure}[H]
  \centering
  \begin{subfigure}[t]{0.45\textwidth}
    \centering
    \includegraphics[width=8cm]{../../../../Image/Book?/8-6e1.png}
    \caption{Kanata trước khi bị chụp.}
  \end{subfigure}
  \quad
  \begin{subfigure}[t]{0.45\textwidth}
    \centering
    \includegraphics[width=8cm]{../../../../Image/Book?/8-6e2.png}
    \caption{Kanata sau khi bị chụp.}
  \end{subfigure}
\end{figure}

Bức tượng ngay sau khi cảm thấy rằng bị những ``con chiên'' trêu đùa, nó lóe mắt đỏ lại, thốt lên: ``\textbf{Chuẩn bị chết đi, Ookami Mio!}''

Nàng Sói đen hoảng hốt: ``Sao lại chỉ có mỗi tôi thôi!? Còn Kuon-senpai\dots\ thì sao!?''

Thấy rằng mình bị đưa vào vô cớ, cậu Tổng trợn mắt: ``Sao lại lôi cả anh vào!?''

Cậu ngước về nàng Dơi, người đang cười ngượng ngùng, rồi cất giọng mắng: ``\textbf{Mà Mel-chan, tự dưng chụp lại em ấy làm cái quái gì!?}''

Vừa nói, cả hai người vừa đung đưa người của nàng Dơi trong sự nóng giận, làm cho cô bị khó thở, chỉ kịp kêu mấy tiếng: ``Chụp hình\dots\ cũng\dots\ vui\dots\ mà!?''

Cậu con trai thét lớn: ``\textbf{Vui cái khỉ mốc ấy!!}''

\img{8-6f}

Cái bức tượng cẩm thạch nói với giọng rất vang: ``\textbf{Ta sẽ đá vào đít các ngươi!}'' và lại lóe sáng một lần nữa.

Cậu Tổng khựng lại, đổ mồ hôi hột với câu nói của bức tượng kia: ``\dots Vãi, đến vị thần Thiên sứ cũng có thể ác miệng được đấy.''

``Thế này là thanh tẩy quá mức rồi!!'' nàng Sói đen phản ứng lại, cố tránh khỏi ánh sáng màu vàng kia.

``Anh còn nghe được điệu \textit{Ode to Joy} nữa cơ\dots''

Bất chợt có giọng hai người từ Thế hệ thứ tư ở phía sau cất lên. Dường như họ đang tám chuyện với nhau một cách vui vẻ. Lamy đang kể một câu chuyện khiến cho Botan khẽ bật cười, không hề để ý gì tới sự việc trước mặt.

\begin{dialogue}
  \speak{Lamy} Bà có biết là Omarun đã làm gì không?
  \speak{Botan} *cười* Làm thế nào cơ?
\end{dialogue}

Mio nhìn thấy họ như vớ được chiếc phao cứu sinh giữa đại dương. ``Em nhìn thấy có Sư tử đi qua kìa!''

Vừa dứt lời, nàng Sư tử trắng liếc về phía bốn người bọn họ.

Không để Botan nói câu gì, Mio ra lệnh: ``\textbf{Mel, chụp đi!}''

Lập tức, nàng Dơi cầm máy ảnh lên, chạy về phía hai người họ, chuẩn bị biến hóa Botan thành một thứ có thể giúp họ.

``Ơ, nhưng mà-'' là những lời Kuon chỉ có thể thốt lên trước khi nàng Sói đen hô: ``\textbf{Làm ngay!}''

``Từ từ đã nào, Botan-'' cậu đang định giải thích thì bị Sói đen cắt lời: ``Cứ chụp đi!''

Mel trước khi ``pô'' một kiểu đã thốt lên nài nỉ: ``Cứu bọn chị với, nữ hoàng của rừng xanh ơi!''

*TÁCH!*

Một thứ ánh sáng lóe lên trước mặt cô và Lamy. Khi ánh sáng dịu đi, trước mặt họ là\dots\ một con mèo mướp nhỏ nhắn.

\begin{figure}[H]
  \centering
  \begin{subfigure}[t]{0.45\textwidth}
    \centering
    \includegraphics[width=8cm]{../../../../Image/Book?/8-6g1.png}
    \caption{Botan trước khi bị chụp.}
  \end{subfigure}
  \quad
  \begin{subfigure}[t]{0.45\textwidth}
    \centering
    \includegraphics[width=8cm]{../../../../Image/Book?/8-6g2.png}
    \caption{Botan sau khi bị chụp.}
  \end{subfigure}
\end{figure}

``\dots anh định nói là Botan nhìn trông hung dữ nhưng mà lại hiền khô mà\dots'' tới lúc này thì đã quá muộn.

``Lại cầm đèn chạy trước ô tô rồi\dots'' Kuon nói trong sự bất lực trước pha xử lý vội vã và đầy bất cẩn của Mio và Mel.

Tới lúc này, gần như họ đã không còn cơ hội để đối phó với bức tượng Thánh thần đang nổi giận kia nữa.

Con mèo ấy cất lên một tiếng ``Meo\~{}'' đầy bùi tai, khiến cho những cô gái mê mệt với những loài sinh vật nhỏ bé này.

Nhìn trước sự dễ thương của chú mèo con, cả Mel và Mio cùng nhau thốt lên: ``Ôi, đáng yêu quá\~{} ♡'' với những khuôn mặt đầy thỏa mãn.

Cậu Tổng thở dài, cam chịu số phận bị bức tượng cẩm thạch kia trừng phạt. Từ phía đằng sau ba người, nó bắt đầu lóe sáng dần, cả thành phố bao phủ một thứ ánh sáng thánh thần đó.

``\textbf{Ta sẽ thanh tẩy các ngươi!!}''

``Ôi thánh thần ơi.''

*VỤT!*

Một luồng ánh sáng rộng lớn bao phủ, thậm chí còn lóa hơn cả đèn nháy từ máy ảnh mà nàng Dơi đã chụp lung tung kia.

\img{8-6h}

Và rồi, khi ánh sáng dần dịu đi\dots

\dots thì bất ngờ thay, chẳng có gì xảy ra cả. Mọi thứ vẫn diễn ra bình thường, mọi người từ trước như thế nào, thì bây giờ vẫn vậy.

Mio và Mel vẫn ngắm con mèo Botan đó một cách âu yếm, thậm chí còn cưng nựng cô ấy đầy thắm thiết. Lamy thì vẫn như hình dạng Yêu tinh lai người như cũ, còn bản thân cậu và A-chan không bị làm sao cả, vẫn còn nguyên vẹn.

Tới lúc này, Kuon mới sực nhớ ra chuyện quan trọng, rồi cũng gật đầu như đã hiểu rõ mọi thứ: ``\textit{Cũng may Kana-chan là thiên sứ học việc nên chưa đủ mạnh để thanh tẩy chúng ta đấy\dots}''

``Giờ, ta phải làm gì với đống hổ lốn này đây? Với cả\dots\ ta có quên ai không nhỉ?''

\dots Câu trả lời là có.

Trong văn phòng \hll, con cừu - vốn là Watame, vẫn đang kêu lên những tiếng be, be inh ỏi. Dường như nó đã thật sự muốn hét lên: ``\textbf{\textit{Mọi người quên mất Watame dễ thương và vô tội này rồi sao!?}}''

Rốt cuộc, sẽ có nhiều thứ để Kuon, Mio và Mel phải giải quyết sau ngày hôm nay.

%==========================================TRUYỆN PHỤ=========================================
\omk

Tất cả mọi người, cùng với Lamy và (mèo) Botan, quay trở về văn phòng. Kuon đang trao điểm với cấp trên về sự việc vừa rồi, trong khi vẫn mắng Mel và Mio.

Nàng Đeo kính thở dài than vãn: ``May cho mấy đứa là không có thương vong nào đấy, chứ không là công ty chúng ta đền ối tiền!''

Cậu Tổng gật gù, ngán ngẩm nhìn nàng Dơi đang cười e thẹn: ``Thế là tạm thế đã. Anh sẽ tính sổ mấy đứa sau. Nhất là em đấy, Mel-chan.''

Cậu nhìn xuống cái máy ảnh kỳ quái trong tay, ánh sáng yếu ớt từ nó vẫn còn đang phát ra chút dư chấn ma thuật. ``Mà\dots\ A-san, chị có biết cái máy ảnh này ở đâu ra không?''

A-chan nhìn chiếc máy ảnh kỳ quặc kia, khoanh tay nhớ lại: ``Hình như là Akai-san và Roboco-san đã mua về từ sạp lễ hội rồi sau đó bỏ quên ở đây.''

Kuon nghe đến đây, đập tay lên trán: ``Cái sạp đó bắn lắm đồ quỷ quái thật\dots\ Mà lại rơi đúng vào hai tên đó nữa\dots\ Họ vẫn chưa chừa từ cái vụ máy ảnh của Quỷ sừng-chan hai năm trước sao?''

Nàng Đeo kính chỉ mỉm cười bất lực: ``Tôi nghĩ chắc họ không cố ý đâu. Thôi, cậu cứ để đấy, ngày mai tôi sẽ tra hỏi hai người đó cho ra nhẽ. Làm phiền cậu rồi, Hikawa-san.''

Cậu Tổng chỉ thở dài chịu trận, gật đầu một cái.

A-chan rời khỏi ăn phòng, nhưng trước khi đi, cô vẫn không quên quay lại nhìn Mel với ánh mắt nghiêm khắc như muốn nhắc nhở: ``Yozora-san, lần sau đừng có tự tiện chụp ảnh người khác mà không xin phép nữa đấy.''

Nàng Dơi chỉ biết cười gượng gạo, ánh mắt len lén nhìn xuống chân như thể muốn chui xuống đất.

Cậu Tổng vẫn cầm chặt cái máy ảnh trong tay, suy nghĩ một lúc lâu. Cậu đưa mắt nhìn quanh phòng: con mèo sư tử vẫn đang chơi đùa với nàng Yêu tinh tuyết, nàng Ma cà rồng nghịch ngợm vẫn đang trêu chọc cừu, mặc cho nó vẫy vùng yếu ớt, còn Sói đen thì đứng đực mặt, ánh mắt đầy khó xử khi nhìn bức tượng cẩm thạch.

\img{8-6i}

Cậu thở dài, cầm lấy cái máy ảnh kỳ quái kia: ``Mấy đứa tránh ra đi, để anh chụp ba cái này.''

Dường như đoán được ý định của cậu, cả nàng Dơi và nàng Yêu tinh tuyết đều tỏ vẻ tiếc nuối.

``Ể\~{}?~? Lamy vẫn đang chơi đùa với con mèo này mà!'' cô nàng ôm con mèo thật chặt, và đối phương cũng không tỏ vẻ  phản đối gì.

``Anh sẽ bảo Okayu-chan thay thế.''

``Còn con cừu này? Em thấy nó cứ loi nhoi trên tay em nãy giờ,'' Mel nghiêng đầu nói, ôm con cừi vẫn đang vùng vẫy trong lòng tay.

``Em quên mất nó là Watame-chan à? Nó muốn được biến trở lại như cũ đấy, Dơi ngố ạ,'' cậu cốc nhẹ vào đầu của nàng Dơi.

``\textbf{\dots Rồi ngươi sẽ được Chúa trời ban phát ân huệ\dots}'' giọng nói từ bức tượng phát ra, vang lên khắp văn phòng, làm mọi người giật mình.

Mio toát mồ hôi hột: ``Dạ, thôi ạ. Tôi cũng không cần. Nhờ anh đó, Kuon-senpai.'' rồi hướng về cậu Tổng.

Ba người nhanh chóng đứng lùi ra sau, để lại ba ``món đồ'' bất đắc dĩ ở chính giữa căn phòng.

Kuon hít sâu một hơi: ``Được rồi, thử xem sao\dots''

Cậu cầm chắc cái máy ảnh, nhắm chuẩn ba mục tiêu rồi bấm nút chụp.

*TÁCH!*

Một luồng sáng trắng chớp lóe lên, chói lòa đến mức mọi người đều phải nheo mắt lại. Không gian bỗng chốc trở nên yên tĩnh một cách kỳ lạ, như thể thời gian vừa bị đóng băng.

Rồi ánh sáng dịu xuống.

Kanata, Watame và Botan ngơ ngác đứng trước mặt họ, rõ ràng là chưa kịp nhận thức chuyện gì vừa xảy ra.

Lamy sau đó vẫn ôm lấy ``người chồng'' của mình đầy thắm thiết, làm cho nàng Sư tử trắng ngơ ngác: ``Hể? Gì vậy? Chuyện gì vừa xảy ra thế?''

Nàng Cừu thì đứng bật dậy, đứng giơ hai tay lên cao vui mừng: ``Yey!! Watame được trở lại rồi này!''

Nàng Thiên sứ thì ngẩn ngơ nhìn mọi người đang đứng xa mình: ``\dots Sao mọi người lại nhìn em như thấy ma thế?''

Khi cả ba món đồ được trở lại bình thường, cậu Tổng mới có thể thở phào nhẹ nhõm bỏ máy ảnh ra.

``Cuối cùng\dots\ cũng xong.''

\begin{center}
  \uline{Ngày hôm sau.}
\end{center}

Khi Roboco và Haato đi làm lại, họ được A-chan và Kuon gọi vào phòng quản lý. Hai người ngồi xuống, vẫn chưa hiểu tại sao họ lại bị gọi như vậy.

``Chuyện gì vậy? Sao hai người nhìn tụi em nghiêm túc thế?'' nàng Song tính hỏi.

Nàng Đeo kính đặt tay lên bàn, nghiêm nghị nói: ``Cái máy ảnh hai đứa mua ở lễ hội\dots\ tụi em có biết nó có phép thuật không?''

Nghe vậy, cả hai cô nàng chớp mắt ngạc nhiên, không thể ngờ tới sự nguy hiểm của nó.

``Ể!? P-Phép thuật ư!?'' / ``\textbf{Cái gì chứ!?}''

Cậu Tổng khoanh tay, nhíu mày lắc đầu ngán ngẩm: ``Nhờ nó mà bọn anh suýt mất ba người đấy.''

Haato và Roboco nhìn nhau, rồi cúi đầu đầy áy náy. Vì đã bày trò nguy hiểm này một lần nữa.

Nàng Người máy bối rối: ``Em cứ tưởng nó chỉ là đồ chơi thôi\dots\'' trong khi nàng Trung học thì vò đầu bứt tai: ``Tại nó nhìn có vẻ thú vị mà!''

Trông thấy hai cô nàng đang hoảng loạn, lại còn giương những đôi mắt tỏ vẻ hối lỗi (mặc dù đây là lần thứ hai còn gặp rắc rối liên quan tới máy ảnh), A-chan cũng chỉ thở dài, không nặng lời thêm nữa. ``Tôi đã quyết định rồi. Cái máy ảnh đó sẽ bị tiêu hủy ngay lập tức. Nhưng mà, vì hai đứa không cố ý nên công ty sẽ đền bù cho hai đứa một khoản tương ứng.''

Cả hai cô nàng khẽ thở phào nhẹ nhõm vì họ không bị phạt. Roboco sáng mắt mừng rỡ, còn Haato thì tươi tỉnh trở lại.

Còn với Mel, người đã gây ra nguyên nhân dẫn tới vụ việc thì\dots

Tại buổi họp cuối ngày của \hll, nàng Ma cà rồng bị nhắc nhở trước toàn bộ thành viên.

Nàng Dơi cúi đầu, mặt đỏ bừng vì xấu hổ. Cô nàng lí nhí: ``Em\dots\ xin lỗi\dots\ Em hứa sẽ không tái phạm nữa\dots''

\dots Cho dù có lẽ đây sẽ không phải lần cuối cô nàng báo hại cả công ty này.

Mio thấy vậy, vỗ vai an ủi tiền bối: ``Nhớ nhé, làm idol là phải giữ hình tượng đấy.''

Sau cùng, mọi chuyện cũng khép lại một cách tốt đẹp (hoặc ít nhất là không có ai phải chịu hậu quả nghiêm trọng hơn).

\dots Nhưng mà, sau vụ này, Kuon chắc chắn sẽ cấm Mel đụng vào bất cứ thiết bị ghi hình nào trong một thời gian dài.

\end{document}